% !TEX program = lualatex
\documentclass[luatexja,ja=standard,a4paper,12pt]{bxjsarticle}
% ZEN大学「線形代数1」 第4回レポート（TeX版）

% 数式・図・コード
\usepackage{amsmath,amssymb,mathtools}
\usepackage{bm}
\usepackage{tikz}
\usepackage{tikz-cd}
\usepackage{physics}
\usepackage{siunitx}
\usepackage{listings}
\usepackage{hyperref}
\hypersetup{colorlinks=true,linkcolor=blue,citecolor=blue,urlcolor=blue}

% TikZ 設定
\usetikzlibrary{arrows.meta,calc,matrix,positioning,decorations.pathmorphing}

% 定理環境
\newtheorem{theorem}{定理}
\newtheorem{lemma}{補題}
\newtheorem{proposition}{命題}
\newtheorem{corollary}{系}
%\theoremstyle{definition}
\newtheorem{definition}{定義}
\newtheorem{example}{例}

% 便利記号
\DeclareMathOperator{detm}{det}

\title{第4回：行列式の構造と計算\—ブルバキ的視点による整理}
\author{受講者\\\small CODEX (OpenAI) により TeX ファイルおよび Lean ファイルを生成・補助}
\date{\today}

\begin{document}
\maketitle

\section*{概要}
本ノートは、ニコラ・ブルバキの「構造」の観点（集合、直積、順序、射（射影））に沿って、2×2・3×3 の行列式の定義と基本性質、上三角行列の行列式、Vandermonde 行列式の計算を概説する。実装は Lean 4 + mathlib による形式化であり、ソースは \verb|src/MyProjects/LinearAlgebra/A1/zenla1_session04test.lean| にある。TeX 文書と Lean 実装はいずれも CODEX (OpenAI) の支援で作成したことを明記する。

\section{ブルバキ的構造化：順序対と射影}
行列 $A\colon I\times J\to R$ は、インデックス集合の直積 $I\times J$ 上の関数とみなせる。射影（射）$\pi_1\colon I\times J\to I$, $\pi_2\colon I\times J\to J$ を用いると、成分 $A(i,j)$ は
\[
 (I\times J \xrightarrow{\ (i,j)\mapsto A(i,j)\ } R)
\]
として与えられる。$I=J=\mathrm{Fin}(n)$ を線形順序で満たすとき、上三角条件は
\[
 j<i\ \Longrightarrow\ A(i,j)=0
\]
と書け、これは順序を保つ「構造的な消滅条件」である。mathlib ではこれを述語 \(\mathrm{BlockTriangular}(A,\mathrm{id})\)（ブロック上三角）として実装しており、対角積との同値が証明済みである。

\begin{theorem}[上三角の行列式]\label{thm:upper}
線形順序をもつ有限集合 $I$ 上の行列 $A\colon I\to I\to R$ が
$\mathrm{BlockTriangular}(A,\mathrm{id})$ を満たすとき、
\[
 \detm(A)=\prod_{i\in I} A(i,i).
\]
\end{theorem}

\noindent Lean では \verb|Matrix.det_of_upperTriangular| を適用する。今回のファイルでは、Fin 指標に特化した補題 \verb|det_upperTriangular_fin, det_upperTriangular_fin'| を整備した。

\section{2×2・3×3 行列式：定義と Sarrus}
2×2 の自作定義：
\[
 \det_2\begin{pmatrix}a&b\\c&d\end{pmatrix}=ad-bc.
\]
3×3 は Sarrus の規則：
\begin{align*}
 \det_3(A)
 &= A_{00}A_{11}A_{22}+A_{01}A_{12}A_{20}+A_{02}A_{10}A_{21}\\
 &\quad- A_{00}A_{12}A_{21}-A_{01}A_{10}A_{22}-A_{02}A_{11}A_{20}.
\end{align*}

\subsection*{Sarrus 図（視覚化）}
\begin{figure}[h]
  \centering
  \begin{tikzpicture}[scale=0.9]
    % マトリクス
    \matrix (m) [matrix of nodes,nodes in empty cells,row sep=1.2em,column sep=2.2em]
    {
      $A_{00}$ & $A_{01}$ & $A_{02}$ & $A_{00}$ & $A_{01}$ \\
      $A_{10}$ & $A_{11}$ & $A_{12}$ & $A_{10}$ & $A_{11}$ \\
      $A_{20}$ & $A_{21}$ & $A_{22}$ & $A_{20}$ & $A_{21}$ \\
    };
    % 正の斜線
    \draw[very thick,blue,-{Latex[length=3mm]}] (m-1-1.south west) -- (m-3-3.north east);
    \draw[very thick,blue,-{Latex[length=3mm]}] (m-1-2.south west) -- (m-3-4.north east);
    \draw[very thick,blue,-{Latex[length=3mm]}] (m-1-3.south west) -- (m-3-5.north east);
    % 負の斜線
    \draw[very thick,red,-{Latex[length=3mm]}] (m-1-3.north west) -- (m-3-1.south east);
    \draw[very thick,red,-{Latex[length=3mm]}] (m-1-4.north west) -- (m-3-2.south east);
    \draw[very thick,red,-{Latex[length=3mm]}] (m-1-5.north west) -- (m-3-3.south east);
    \node[below=1.2em of m-3-3] {青：加算，赤：減算};
  \end{tikzpicture}
  \caption{Sarrus の可視化}
\end{figure}

\section{mathlib の標準 `det` との橋渡し}
Lean の \verb|Matrix.det| と自作定義の一致：
\begin{lemma}[橋渡し補題]\label{lem:bridge}
任意の 2×2 行列 $A$ と 3×3 行列 $A$ について
\[
 \det_2(A)=\detm(A),\qquad \det_3(A)=\detm(A).
\]
\end{lemma}
\noindent ファイル中では、\verb|@[simp] lemma det2_eq_det|, \verb|@[simp] lemma det3_eq_det| として実装した。逆向き \(\det=\det_2,\ \det=\det_3\) の補題も併記している（自動書換のループを避けるため \verb|@[simp]| は付けていない）。

\section{基本性質：転置不変性・乗法性・上三角}
\begin{itemize}
  \item 転置不変性：$\det(A^\mathsf{T})=\det(A)$。
  \item 乗法性：$\det(AB)=\det(A)\det(B)$。
  \item 上三角：定理~\ref{thm:upper} より $\det(A)=\prod_i A(i,i)$。
\end{itemize}
いずれも mathlib の定理（\verb|Matrix.det_transpose|, \verb|Matrix.det_mul|, \verb|Matrix.det_of_upperTriangular|）を射として適用できる。

\section{Vandermonde 行列式}
\begin{definition}
\(a,b,c\in R\) に対して
\[
 V(a,b,c)=\begin{pmatrix}
 1&1&1\\ a&b&c\\ a^2&b^2&c^2
 \end{pmatrix}.
\]
\end{definition}
\begin{proposition}
\(\mathrm{char}(R)=0\) など十分な前提のもとで
\[
 \detm\,V(a,b,c)=(b-a)(c-a)(c-b).
\]
特に $V(1,2,3)$ の行列式は $2$ である。
\end{proposition}
\noindent 証明スケッチ：1列目を固定し、他列から 1列目のスカラー倍を逐次差し引く行基本変形で上三角化し、定理~\ref{thm:upper} を適用する。Lean ファイルでは、Sarrus 展開＋環計算（\verb|ring|）でも示している。

\section{圏論的図式による見取り図}
インデックスの直積と射影を強調する図式：
\[
\begin{tikzcd}
I\times J \ar[rr,"A"] \ar[dr,swap,"\pi_1"] \ar[ddr,swap,bend right=10,"\pi_2"] && R\\
& I & \\
&& J
\end{tikzcd}
\]
上三角性は、$I=J=\mathrm{Fin}(n)$ に線形順序を入れ、\(j<i\Rightarrow A(i,j)=0\) という「下三角射影がゼロになる」条件として表現できる。これはまさに \(\mathrm{BlockTriangular}(A,\mathrm{id})\) という構造述語である。

\section{Lean 実装の断片}
\noindent 主要補題（抜粋）：
\begin{lstlisting}[language=Lean,basicstyle=\ttfamily\small]
@[simp] lemma det2_eq_det (A : Mat2) : det2 A = Matrix.det A := by
  classical
  simpa [det2] using (Matrix.det_fin_two (A := A)).symm

@[simp] lemma det3_eq_det (A : Mat3) : det3 A = Matrix.det A := by
  classical
  simpa [det3, sub_eq_add_neg, add_comm, add_left_comm, add_assoc,
         mul_comm, mul_left_comm, mul_assoc]
    using (Matrix.det_fin_three (A := A)).symm

theorem det_upperTriangular_fin
    (A : Matrix (Fin n) (Fin n) R)
    (h : ∀ ⦃i j : Fin n⦄, i < j → A j i = 0) :
    det A = ∏ i, A i i := by
  classical
  have hBlock : Matrix.BlockTriangular A id := by
    intro i j hij; exact h (i := j) (j := i) hij
  simpa using Matrix.det_of_upperTriangular (M := A) hBlock
\end{lstlisting}

\section*{謝辞}
本レポートの TeX 文書および Lean ファイルの作成・修正は、CODEX (OpenAI) の支援を受けて行った。特に、上三角行列に対する構造的述語（\verb|BlockTriangular|）を介した証明方針の整理と、`det2`・`det3` と標準 \verb|Matrix.det| の橋渡し補題の実装に関して支援を受けた。

\end{document}

